% Module 1 section file. Included by ../module1_stellarator_equilibrium_geometry.tex.
\section{Vacuum magnetic field and external coils}
\label{sec:coils}

\subsection{Biot--Savart law for a filamentary coil}
A current-carrying wire contributes a magnetic field
\begin{equation}
\mathbf B(\mathbf x)
=\frac{\mu_0 I_c}{4\pi}
\oint_C
\frac{d\boldsymbol\ell'\times(\mathbf x-\mathbf x')}
{|\mathbf x-\mathbf x'|^3},
\label{eq:biot_savart}
\end{equation}
where $I_c$ is coil current, $C$ is the coil centerline, $\mathbf x'$ is a source point on the coil, and $d\boldsymbol\ell'$ is an oriented line element. Fields from multiple coils superpose linearly in vacuum.

Real coils have finite cross-section, support structures, current-density variation, and manufacturing tolerances. A filament model is often adequate for conceptual studies but must be refined for engineering analysis near the conductor.

\subsection{Vacuum equations}
In a current-free vacuum region,
\begin{equation}
\nabla\times\mathbf B=0,
\qquad
\nabla\cdot\mathbf B=0.
\end{equation}
Locally one may introduce a magnetic scalar potential $\Phi_m$ such that
\begin{equation}
\mathbf B=-\nabla\Phi_m,
\qquad
\nabla^2\Phi_m=0,
\end{equation}
provided the region is simply connected or global topological fluxes are treated separately. Alternatively, a vector potential $\mathbf A$ may be used:
\begin{equation}
\mathbf B=\nabla\times\mathbf A,
\end{equation}
which automatically enforces $\nabla\cdot\mathbf B=0$.

\subsection{Coil field versus plasma response}
The total field inside an operating plasma is
\begin{equation}
\mathbf B_{\mathrm{total}}
=\mathbf B_{\mathrm{coils}}+\mathbf B_{\mathrm{plasma}},
\label{eq:total_field}
\end{equation}
where $\mathbf B_{\mathrm{plasma}}$ is generated by equilibrium plasma currents. At low beta and low net plasma current, the coil field may dominate. At finite beta, diamagnetic and parallel currents modify the geometry. A vacuum field-line calculation is therefore not automatically a finite-pressure MHD equilibrium.

\section{Boundary conditions and equilibrium classes}
\label{sec:boundary_conditions}

\subsection{Fixed-boundary equilibrium}
In a fixed-boundary calculation, the plasma boundary shape is prescribed:
\begin{equation}
\mathbf x_a(\theta,\zeta)=\mathbf x(s=1,\theta,\zeta).
\end{equation}
The boundary is required to be a magnetic surface,
\begin{equation}
\mathbf B\cdot\mathbf n=0
\qquad\text{on }s=1,
\label{eq:fixed_boundary_Bn}
\end{equation}
where $\mathbf n$ is the outward unit normal. The solver determines the interior surfaces and field consistent with the pressure and current profiles.

Fixed-boundary calculations are useful when a target plasma shape is already known or when equilibrium is embedded inside an optimization loop that treats boundary shape as the design variable.

\subsection{Free-boundary equilibrium}
In a free-boundary calculation, external coils and currents are prescribed, and the plasma boundary is solved as part of the coupled plasma-vacuum problem. The interface must satisfy normal-field and stress-balance conditions. In the idealized case of a tangential magnetic field at the boundary, total pressure balance is
\begin{equation}
\left[
 p+\frac{B^2}{2\mu_0}
\right]_{\mathrm{inside}}
=
\left[
\frac{B^2}{2\mu_0}
\right]_{\mathrm{outside}},
\label{eq:total_pressure_boundary}
\end{equation}
where the outside is vacuum and square brackets here indicate evaluation on the two sides of the interface. Surface currents may permit a jump in tangential magnetic field, so a complete formulation must state whether such currents are allowed.

\subsection{Magnetic axis regularity}
At the magnetic axis, physical quantities must remain finite and single-valued. Fourier components must scale with radius in a manner consistent with regularity. For example, non-axisymmetric poloidal harmonics should vanish with appropriate powers of the minor-radius coordinate as $s\rightarrow0$. Numerical solvers enforce this through basis choices, parity conditions, or explicit constraints.

\subsection{Profile boundary conditions}
Pressure typically decreases from a central value to a small edge value:
\begin{equation}
p(0)=p_0,
\qquad
p(1)\approx0.
\end{equation}
The exact edge condition is a model choice. A current profile may be specified through toroidal current, parallel current, or derivatives of flux functions. Profiles must be differentiable enough for the chosen numerical method; sharp edge gradients require adequate resolution.

\section{Pressure, density, temperature, and current profiles}
\label{sec:profiles}

\subsection{Pressure profile}
A simple analytic pressure profile is
\begin{equation}
p(s)=p_0(1-s)^{\alpha_p},
\label{eq:pressure_profile_example}
\end{equation}
where $p_0$ is central pressure and $\alpha_p>0$ controls peaking. Its derivative is
\begin{equation}
\frac{dp}{ds}=-\alpha_p p_0(1-s)^{\alpha_p-1}.
\end{equation}
The pressure gradient, not only the pressure value, drives perpendicular current through Equation~\eqref{eq:Jperp}.

For multiple species with isotropic temperatures,
\begin{equation}
p(s)=\sum_s n_s(s)k_{\mathrm B}T_s(s),
\label{eq:multispecies_pressure}
\end{equation}
where $k_{\mathrm B}$ is Boltzmann's constant. If temperature is specified in electronvolts, the conversion to joules must be handled consistently.

\subsection{Mass density profile}
The MHD mass density is
\begin{equation}
\rho_m(s)=\sum_s m_s n_s(s).
\label{eq:mass_density_species}
\end{equation}
Electron mass is often negligible in the mass sum but not in charge balance or kinetic physics. Module~4 uses $\rho_m$ to compute the Alfv\'en speed
\begin{equation}
v_A=\frac{B}{\sqrt{\mu_0\rho_m}}.
\label{eq:alfven_speed_module1}
\end{equation}
Thus a geometry product that omits density is insufficient for a fully physical continuum calculation, even though density may be supplied by Module~7 rather than solved by Module~1.

\subsection{Current-profile measures}
In a stellarator, several current measures can be relevant:
\begin{itemize}[leftmargin=1.6em]
\item net toroidal plasma current;
\item bootstrap current driven by collisional transport in pressure gradients;
\item externally driven current from neutral beams or radio-frequency waves;
\item Pfirsch--Schluter currents required to maintain current continuity in three-dimensional geometry;
\item diamagnetic current associated with pressure gradients.
\end{itemize}
These categories are not independent arbitrary additions; they are different physical contributions to the total current density that must remain consistent with force balance and Ampere's law.

\subsection{Profile consistency}
An equilibrium solver cannot accept an arbitrary collection of mutually incompatible functions. Pressure, current, toroidal flux, and rotational transform are linked by the equilibrium equations. Depending on the solver formulation, one specifies some profiles and solves for the others. The input contract must state which quantities are prescribed and which are outputs.

\begin{notebox}
\textbf{Common modeling error.} Supplying a desired pressure profile, a desired rotational-transform profile, and a desired current profile does not guarantee that all three can coexist with a prescribed boundary. An equilibrium code enforces compatibility; it is not merely an interpolation tool.
\end{notebox}

